\section{Conclusão}
\label{sec:conclusao}

Apresentamos o DRM Transformer, um Transformer decoder-only que substitui a
atenção por produto escalar por uma atenção geodésica em um Directional
Relational Manifold aprendido. A arquitetura combina MetricNet low-rank,
deformação gravitacional, gamma-scaling relativístico e dimensionalidade
variável. A partir dos papers DRM e de dinâmica relativística, implementamos
uma versão neural treinável da ideia de dimensionalidade emergente e estrutura
toroidal.

O resultado mais importante até o momento é a observação de uma assinatura
homológica compatível com uma superfície fechada quase-toroidal e robusta a
sementes no baseline small\_3.5M. Os controles sugerem que o treinamento
desloca o manifold DRM de uma topologia inicial de gênero alto e instável para
essa assinatura: em duas de três sementes treinadas, a topologia mediana
coincide com $H_1=2$, $H_2=1$, assinatura esperada de $T^2$; em todas as
sementes treinadas há sinal \emph{Near T2}, enquanto a inicialização aleatória
não apresenta esse sinal. Esse resultado não encerra a validação do método,
pois o critério estrito $f_{T^2}\geq 0{,}60$ não foi atingido, mas demonstra que
a geometria DRM pode ser operacionalizada em modelos de linguagem e avaliada
por ferramentas topológicas objetivas.

Trabalhos futuros incluem escalar o modelo, testar ablações em 50M, 350M e
1B+ parâmetros, comparar contra atenção padrão sob orçamento equivalente e
estudar a relação entre topologia DRM e perplexidade. Um ponto central de
validação será comparar quatro regimes: toro explícito, toro com annealing até
zero, geometria genérica sem alvo toroidal e ausência de regularização
toroidal. Se regimes sem alvo toroidal mantiverem $H_1=2$, $H_2=1$ com
estabilidade maior em escala, a evidência de emergência topológica será
substancialmente mais forte do que no regime explicitamente prescritivo.
